\documentclass[11pt]{article}
\usepackage{amsmath,amssymb}

\begin{document}
% Standalone preview of manuscript Section 3.3. In the full manuscript,
% copy the content from \subsection through the last paragraph; the journal
% class will number the equations automatically.
\setcounter{section}{3}
\setcounter{subsection}{2}

\subsection{Hyperspectral preprocessing}

\subsubsection{Cube Canonicalisation and Leaf Masking}
The TIFF cubes were loaded in the order height $\times$ width $\times$ band
($48 \times 352 \times 204$). When required, band-first arrays were reordered
to this convention, and no spatial resizing was applied to the global cubes at this stage. The stored unsigned
16-bit intensity values were scaled to the range $[0,1]$ as defined in
Eq.~\eqref{eq:normalization}:

\begin{equation}
U_i(p,b)=\frac{X_i(p,b)}{65{,}535},
\label{eq:normalization}
\end{equation}

where $X_i$ denotes the stored hyperspectral cube of image $i$, $U_i$ represents
the corresponding unit-scaled cube, $p$ indicates a spatial pixel location, and
$b$ denotes the spectral band index. A deterministic foreground mask was then
generated directly from the spectral measurements. A pixel was assigned to the
foreground when at least 75\% of its 204 spectral bands contained non-zero
intensity values, as specified in Eq.~\eqref{eq:foreground_mask}:

\begin{equation}
M_i^{(0)}(p)=
\mathbf{1}
\left[
\frac{1}{204}
\sum_{b=1}^{204}
\mathbf{1}\{U_i(p,b)>0\}
\geq 0.75
\right],
\label{eq:foreground_mask}
\end{equation}

where $M_i^{(0)}$ represents the initial mask before post-processing. Connected
components containing fewer than 20 pixels were removed, and holes within the
remaining components were subsequently filled. Pixels outside the final leaf
mask were assigned zero intensity. The mask was derived solely from the hyperspectral measurements and did not use prediction targets or
manually annotated segmentation masks.

\subsubsection{Fold-specific normalisation}
For each spectral band, the mean and standard deviation were calculated only
from the pooled foreground pixels in the 1,272 outer-training images. Each
foreground pixel was normalised by subtracting the corresponding band mean and
dividing by its standard deviation for that fold. Background pixels remained
zero after normalisation. The fold-specific statistics and transformation are
defined in Eq.~\eqref{eq:spectral_standardization}.

\begin{equation}
\begin{aligned}
N_f
&=
\sum_{i\in\mathcal{T}_f}\sum_p M_i(p),
\\[4pt]
\mu_{f,b}
&=
\frac{1}{N_f}
\sum_{i\in\mathcal{T}_f}\sum_p
M_i(p)U_i(p,b),
\\[4pt]
\sigma_{f,b}
&=
\sqrt{
\max\left(
\frac{1}{N_f}
\sum_{i\in\mathcal{T}_f}\sum_p
M_i(p)
\left(U_i(p,b)-\mu_{f,b}\right)^2,
10^{-12}
\right)
},
\\[4pt]
\widetilde{U}_{i,f}(p,b)
&=
M_i(p)
\frac{U_i(p,b)-\mu_{f,b}}
{\sigma_{f,b}} .
\end{aligned}
\label{eq:spectral_standardization}
\end{equation}

Here, $\mathcal{T}_f$ is the outer-training set for fold $f$, $M_i$ is the
final leaf mask, and $N_f$ is the total number of foreground pixels used to
estimate the band statistics for that fold. The three continuous targets were
also standardised separately using the means and population standard deviations
calculated from the same outer-training partition. Before evaluation and
reporting, the predicted values were transformed back to their original
scales, as shown in Eq.~\eqref{eq:target_standardization}.

\begin{equation}
\begin{aligned}
n_f
&=
|\mathcal{T}_f|
=
1{,}272,
\\[4pt]
\mu_{f,t}
&=
\frac{1}{n_f}
\sum_{i\in\mathcal{T}_f}
y_{i,t},
\\[4pt]
s_{f,t}
&=
\sqrt{
\frac{1}{n_f}
\sum_{i\in\mathcal{T}_f}
\left(y_{i,t}-\mu_{f,t}\right)^2
},
\\[4pt]
z_{i,t}
&=
\frac{y_{i,t}-\mu_{f,t}}
{s_{f,t}},
\qquad
\widehat{y}_{i,t}
=
\mu_{f,t}
+
s_{f,t}\widehat{z}_{i,t}.
\end{aligned}
\label{eq:target_standardization}
\end{equation}

Here, $t$ denotes one of the three continuous targets, and
$\widehat{z}_{i,t}$ represents its prediction on the standardised scale.

\subsubsection{Training augmentation}
Data augmentation was performed exclusively during PhytoCore training, following
a fixed sequence: multiplicative intensity gain, additive Gaussian noise,
smooth spectral jitter, and spatial translation. Intensity gain was applied
with a probability of 0.5, using a factor sampled from the range
$[0.98,1.02]$. If more than 5\% of foreground spectral values were already
saturated at 1.0, the upper bound was reduced to 1.0. Gaussian noise was
applied with probability 0.5, with its standard deviation sampled from
$[0,0.002]$. Spectral jitter was applied with probability 0.35. A random
band-wise curve was smoothed with a Gaussian filter ($\sigma=8$ bands),
centred and standardised, then scaled by an amplitude sampled from
$[0,0.01]$. The resulting multiplicative factors were clipped to
$[0.95,1.05]$. Translation was applied with probability 0.5, with maximum
displacements of eight pixels horizontally and two pixels vertically;
allowable shifts were further limited to keep the leaf within the image.
Intensities were clipped to $[0,1]$, and background pixels were reset to zero
before fold-specific normalisation.

\end{document}
